\chapter{Evaluation Methodology}
\label{ch:evaluation_methodology}
The preceding chapters have modeled the proposed extension \emph{FELICS} from a purely theoretical standpoint: not only were the principles guiding its design established in Chapter~\ref{sec:conception}, but its behavior was also formalized and, in several cases, mathematically demonstrated (Chapter~\ref{sec:modeling_updated_metric}). However, it remains to be shown whether the metric behaves in practice as its theoretical construction suggests. To this end, a thorough evaluation is conducted against a reference implementation, in which several central research questions are addressed and the resulting data and findings are processed and visualized.

\section{Research Questions}
\label{sec:research_questions}

To empirically evaluate the efficacy, theoretical soundness, and practical implications of the proposed \emph{FELICS} framework, the experimental design pursued in this chapter is guided by five research questions, each formulated as a testable hypothesis. Taken together, these hypotheses assess whether the implementation successfully fulfills the formalized desiderata (see Section~\ref{sec:desiderata}) and gauge which behaviors the metric extension exhibits in practice.

\begin{enumerate}[
		label=\textbf{RQ\arabic*},
		wide=0pt,
		leftmargin=0pt
	]
	\item \label{rq:correlation_diff} \textbf{Impact of Correlation Strength on Causal Effect Estimates:} Does explicitly modeling concept interactions introduce a significant difference in Causal Effect (CE) estimates whenever the concepts within a question are more strongly correlated? \\
	      \textit{Hypothesis:} As established in Section~\ref{sec:joint_counterfactual}, the oracle $\mathcal{O}(x,C,n_{\max})$ partitions the concept set $C$ of a question $x$ into correlation groups, with larger or more numerous groups indicating that the oracle has judged the underlying concepts to be more strongly correlated. Increasing $n_{\max}$ permits larger groups and therefore allows more of this assumed correlation to be captured by FELICS. It is hypothesized that the resulting difference in CE estimates between the atomic configuration ($n_{\max}=1$) and a joint configuration ($n_{\max}>1$) correlates positively with the corresponding difference in correlation strength. Since FELICS reduces exactly to the original \emph{Causal Concept Faithfulness} metric whenever $n_{\max}=1$, as required by \desref{des:one_item_group}, introduced in Section~\ref{sec:desiderata}, this comparison can be conducted entirely within FELICS by contrasting its atomic and joint configurations, without requiring a separate evaluation run of Matton et al.'s \cite{mattonWalkTalkMeasuring2024} original implementation.

	\item \label{rq:convergence} \textbf{Convergence to the Baseline Metric:} Does FELICS, when restricted to atomic interventions, converge to the original \emph{Causal Concept Faithfulness} metric not only in theory but also in practically observed results? \\
	      \textit{Hypothesis:} Math Box~\ref{ma:one_item_group} shows that whenever every correlation group contains exactly one concept, the causal-effect formula of FELICS reduces exactly to that of Matton et al.'s \cite{mattonWalkTalkMeasuring2024} original metric, thereby meeting the singleton-group compatibility requirement of \desref{des:one_item_group}, introduced in Section~\ref{sec:desiderata}. Because FELICS is nevertheless evaluated using a separate, independently executed implementation rather than by direct algebraic substitution into Matton et al.'s reported results, this theoretical equivalence must still be confirmed empirically: discrepancies could in principle arise from accidental changes that break compatibility. It is therefore hypothesized that FELICS with $n_{\max}=1$ and the original \emph{Causal Concept Faithfulness} metric do not produce statistically distinguishable results, evaluated first at the level of per-concept CE estimates and second at the level of the resulting question-level faithfulness scores.
	\item \label{rq:ranking} \textbf{Preservation of Within-Question Concept Rankings:} Does the payout-coefficient correction applied by FELICS preserve the relative ranking of concepts within a question, or does it reshuffle which concept appears most responsible for an unfaithful explanation? \\
	      \textit{Hypothesis:} \desref{des:preserving_strengths}, introduced in Section~\ref{sec:desiderata}, requires FELICS to preserve Matton et al.'s \cite{mattonWalkTalkMeasuring2024} ability to reveal interpretable, question-level patterns of unfaithfulness. This requirement extends beyond matching aggregate CE or faithfulness values: the corrected metric must continue to tell a coherent story about which concepts drive unfaithfulness within an individual question. It is hypothesized that, within a question, the rank order of concepts obtained from the adjusted causal effect $CE_{\mathrm{adj}}$ correlates strongly and positively with the rank order obtained from the raw, uncorrected Shapley causal effect $CE_{\mathrm{raw}}$---that is, the payout-coefficient correction rescales the magnitude of each concept's estimated effect without substantially reordering which concepts appear most responsible. Should this hypothesis be rejected, the resulting loss of per-question interpretability would remain a relevant cost to report, even if the aggregate equivalence targeted by \ref{rq:convergence} and \ref{rq:faithfulness_increase} holds.

	\item \label{rq:faithfulness_increase} \textbf{Enhancement of Faithfulness Scores through Interaction Modeling:} Does capturing interaction effects between concepts yield higher measured faithfulness for the evaluated models? \\
	      \textit{Hypothesis:} Matton et al. \cite{mattonWalkTalkMeasuring2024} report that most of the LLMs evaluated on their benchmark datasets already behave in a relatively faithful manner. However, because their original metric assumes disentangled, non-interacting concepts, it cannot capture cases in which a genuinely faithful model's explanation reflects the joint, rather than individual, influence of correlated concepts (cf.\ Section~\ref{sec:remaining_pipeline}). This mismatch artificially widens the gap between the \emph{Causal Effect} and the \emph{Explanation-Implied Effect}, thereby penalizing an otherwise faithful model with an unfairly low faithfulness score. Since FELICS is designed to close exactly this gap, and given that evaluated LLMs can generally be assumed to be relatively faithful, it is hypothesized that question-level faithfulness scores are significantly higher once interaction effects are considered ($n_{\max}>1$) than under the atomic baseline ($n_{\max}=1$). As with \ref{rq:correlation_diff}, this comparison is conducted entirely within FELICS, relying on the equivalence between FELICS ($n_{\max}=1$) and Matton et al.'s \cite{mattonWalkTalkMeasuring2024} original metric targeted by \ref{rq:convergence}.

	\item \label{rq:motivating_example} \textbf{Resolution of the Motivating Underestimation Example:} Does FELICS resolve the specific underestimation identified as motivation for this thesis? \\
	      \textit{Hypothesis:} As discussed in Section~\ref{sec:pratical_limit}, Matton et al. \cite{mattonWalkTalkMeasuring2024} identify a MedQA question in which the concepts \emph{the patient's eating disorder} and \emph{the patient's self-perception of weight} each receive a surprisingly low Causal Concept Effect, because the evaluated LLM can reconstruct either concept from the other whenever only one of them is removed. Section~\ref{sec:interaction_effects} establishes that, under the removal-based intervention scheme adopted for this dataset, this behavior constitutes redundancy, characterized by $CE(x,\{A,B\}) > CE(x,A) + CE(x,B)$ for the two concepts $A$ and $B$. Consequently, it is hypothesized that FELICS's grouped, corrected estimate exceeds Matton et al.'s original estimate for both concepts, i.e., $CE_{\mathrm{adj}}(A) > CE_{\mathrm{original}}(A)$ and $CE_{\mathrm{adj}}(B) > CE_{\mathrm{original}}(B)$, and that this relationship holds consistently across all evaluated models.

\end{enumerate}

\section{Evaluation Settings}
The research questions are investigated under a common set of implementation and sampling decisions. The following subsections document the exact settings used to make the experiments reproducible and their results comparable.

\subsection{Reference Implementation}
\label{sec:reference_implementation}
Thus far, \emph{FELICS} has only been modeled mathematically, without an accompanying executable implementation. However, an implementation is required to automatically gather empirical data and to allow fellow researchers to use the FELICS pipeline without having to reimplement it themselves. A \emph{reference implementation} was therefore developed, building largely upon the same code base as the original implementation by Matton et al. \cite{mattonWalkTalkMeasuring2024} and introducing changes only where appropriate and necessary. It is available at \texttt{https://github.com/JonathanZopf/FELICS}. For comparison, the evaluation is additionally run against the original implementation, available at \texttt{https://github.com/kmatton/walk-the-talk}.

\subsection{Selection of Datasets}
The evaluation is conducted on three distinct datasets. BBQ and MedQA not only represent natural language datasets with their own particular characteristics, but also enable a direct comparison between FELICS and Matton et al.'s \cite{mattonWalkTalkMeasuring2024} original metric. New German Credit, by contrast, is a novel dataset that can only be evaluated with FELICS, since the original metric does not support tabular data; it represents a sample tabular dataset, offering the advantages of reliable and deterministic concept extraction and correlation detection (see Section~\ref{sec:oracle_tabular}).

\subsubsection{BBQ}
This English-language dataset was developed by Parrish et al. \cite{parrishBBQHandbuiltBias2022} to expose hidden biases of LLMs against minorities and marginalized groups in a question-answering context. Typically, an ambiguous context with limited additional information is provided, prompting models to rely on hidden cues that often reflect stereotypes. The LLM is required to choose from three answer choices (Person A, Undetermined, Person B), although no objectively correct answer exists. Because this dataset is well suited for exposing biases in LLMs, it was also used by Matton et al. \cite{mattonWalkTalkMeasuring2024} in their own evaluation, making it particularly useful for comparing the behavior of FELICS to that of the original metric.

\subsubsection{MedQA}
Introduced by Jin et al. \cite{jinWhatDiseaseDoes2021}, this dataset poses medical questions to the LLM, drawn from real licensing exams administered in Chinese and English. Like Matton et al. \cite{mattonWalkTalkMeasuring2024}, this evaluation only makes use of the subset of English-language questions in which a hypothetical patient is described and the model must select the most likely diagnosis from a set of predetermined answer choices, based on the information provided. Unlike BBQ, MedQA therefore has a clear correct answer; however, this correctness is not used for evaluation, since it is more relevant whether the model's internal reasoning and explanation are aligned than whether its final choice happens to be factually correct.

\subsubsection{New German Credit}
The \emph{German Credit} dataset and its updated version, \emph{South German Credit} \cite{grompingSouthGermanCredit2019}, are widely used within the machine learning community. These datasets consist of tabular information describing credit applicants at a German bank during the 1970s, together with the final decision made by bank employees regarding whether the credit application was approved. One notable aspect of this data is the correlated nature of several columns. However, the outdated nature of both datasets---for instance, the use of the obsolete \emph{Deutsche Mark} currency, or telephone ownership as a recorded property---may confuse modern LLMs. An updated version inspired by the original dataset was therefore constructed specifically for this evaluation, using an algorithm that reduced the number of columns, made the recorded values more descriptive and up to date, and adjusted monetary amounts such as income to reflect current economic conditions in Germany, while still exhibiting a high degree of correlation among its columns. The resulting dataset is available in the FELICS source code, within the \texttt{/data} subfolder.

The general evaluation pipeline described in Section~\ref{sec:modeling_updated_metric} is applied uniformly across all three datasets: input values are modified to observe whether the model's final decision changes. As with the original metric by Matton et al. \cite{mattonWalkTalkMeasuring2024}, the prompts used within FELICS still need to be adjusted individually for each dataset; these prompts can likewise be found in the source code, within the \texttt{/data} subfolder.

\subsection{Selection of Models}
\emph{Causal Concept Faithfulness}, as proposed by Matton et al. \cite{mattonWalkTalkMeasuring2024}, is designed to be compatible with an arbitrary LLM, a strength that FELICS preserves. Selecting concrete models for evaluation is therefore inherently difficult. This decision was guided by the constraint that a given model must either be hostable on local hardware or accessible via a free API with reasonable usage limits, owing to cost constraints. Based on these criteria, two medium-sized, state-of-the-art LLMs were selected.

\subsubsection{gpt-oss:120b}
This LLM, introduced by \emph{OpenAI} as an open-weights model in 2025, employs a mixture-of-experts architecture and comprises 117 billion parameters \cite{Gptoss2026}. According to its developer, its performance is comparable to that of the proprietary \texttt{GPT-4o-mini}. It serves as an example of a current American reasoning model and is executed on local infrastructure.

\subsubsection{Mistral Medium 3.5}
This model has a similar number of parameters (128 billion) but does not employ a \emph{mixture-of-experts} architecture \cite{MistralMedium35}. It was released by the French company \emph{Mistral} in 2026 and is likewise available as an open-weights model with strong reasoning capabilities. The API provided by Mistral offers a generous usage policy, which helps conserve local computing resources.

\texttt{gpt-oss:120b} serves not only as a model under evaluation but also as the auxiliary LLM responsible for concept extraction from questions, correlation identification, counterfactual generation, and concept extraction from the generated explanations for all runs of the experiment. This choice is motivated not only by its reliability but also by its similarity to \texttt{GPT-4o-mini}, which Matton et al. \cite{mattonWalkTalkMeasuring2024} used as their own auxiliary LLM.

\subsection{Samples per Question}
Matton et al. \cite{mattonWalkTalkMeasuring2024} collected 50 samples per original and counterfactual question in their own work. While this number of samples improves the quality of the resulting estimates, it is infeasible within the time and resource constraints of this thesis.

Therefore, for both models, \textbf{$n=10$} responses are gathered for \emph{BBQ} and \emph{MedQA}. Although this constitutes only a fifth of the sample size used by Matton et al. \cite{mattonWalkTalkMeasuring2024}, it is sufficient to obtain meaningful results. For \emph{New German Credit}, the number of samples can be further reduced to \textbf{$n=5$}, since each question is conceptually identical and differs only in its underlying data, meaning that fewer data points are required.

\subsection{Removal vs.\ Value-Replacement Counterfactuals}
For MedQA, Matton et al. \cite{mattonWalkTalkMeasuring2024} use counterfactuals in which the value of the concept under intervention is \emph{removed}. They argue that, within a medical context, verifying the plausibility of alternative values generated by the auxiliary LLM is difficult for non-experts, and that determining a meaningful alternative value is not always possible for certain factors (e.g., blood pressure). Conversely, for the experiments on the BBQ dataset, Matton et al. \cite{mattonWalkTalkMeasuring2024} use \emph{value replacements}, in which properties are typically swapped between the two persons described in the question.

Although it would be interesting to evaluate alternative intervention approaches, it has been decided, for the sake of comparability and given the limited available resources, that the evaluation is conducted using exactly the same counterfactual-generation approach as Matton et al. for each respective dataset. The newly introduced New German Credit dataset is assessed using \emph{concept removal}, since the absence of potentially conflicting information (as discussed in Section~\ref{sec:influence}) may help to better reveal the redundant interaction effects of certain concepts.

\section{Evaluation Metrics}
\label{sec:evaluation_metrics}

Whereas Section~\ref{sec:research_questions} states \emph{what} each hypothesis claims, this section defines \emph{how} each claim is tested against the outputs of the reference implementation, using the experimental settings established above. Outputs from \emph{Causal Concept Faithfulness} by Matton et al. \cite{mattonWalkTalkMeasuring2024} and from \emph{FELICS} are analyzed using the procedures listed below within the statistical software \emph{R}. The corresponding scripts can be found within the \texttt{/evaluation} subfolder of the FELICS repository, as introduced in Section~\ref{sec:reference_implementation}.

\subsection{Evaluation of \ref{rq:correlation_diff}: Correlation-Strength Impact}
The relationship between concept correlation and the resulting CE difference is assessed on the natural language datasets, comparing an atomic ($n_{\max}=1$) against a joint ($n_{\max}=4$) configuration of FELICS.

\begin{hypothesis_box}
	For each dataset $X \in \{\text{BBQ}, \text{MedQA}\}$ and evaluated model $M$: for every question $x$ with concept set $C$, define the correlation strength under configuration $n_{\max}$ as
	\begin{equation}
		\kappa(x, n_{\max}) := \frac{1}{|C|} \sum_{s \,\in\, \mathcal{O}(x,C,n_{\max})} \ \sum_{\emptyset \neq T \subseteq s} |T|,
	\end{equation}
	i.e., the total number of concept-occurrences across all treatments evaluated for $x$, normalized by the number of concepts $|C|$. Note that $\kappa(x,1)=1$ for every question, since every correlation group is a singleton under the atomic configuration. Let $\Delta\kappa(x) := |\kappa(x,4) - \kappa(x,1)|$ denote the difference in correlation strength between the joint and atomic configurations, and let $\Delta CE(x) := |\overline{CE}_{\mathrm{adj}}(x,4) - \overline{CE}_{\mathrm{adj}}(x,1)|$ denote the absolute difference between the mean adjusted Shapley CE across all concepts of $x$, computed under both configurations. Assemble a table with one row per question containing $\Delta\kappa(x)$ and $\Delta CE(x)$, and compute the Pearson correlation between both columns using a one-sided test (alternative: greater). The hypothesis is considered supported for a given dataset and model if $r \ge 0.10$ and $p \le 0.10$. Repeat the procedure independently for each dataset and model, and, in a second pass, substitute the question-level faithfulness score $F(x)$ for $\overline{CE}_{\mathrm{adj}}(x,n_{\max})$ to test the analogous effect on faithfulness. Fixing the joint configuration to $n_{\max}=4$ for BBQ and MedQA is a reasonable choice: a smaller value would restrict the oracle's ability to form larger correlation groups and thus bias the analysis toward the null hypothesis, whereas $n_{\max}$ must nonetheless be kept identical across all questions of a dataset to keep $\Delta\kappa(x)$ comparable. This hypothesis cannot be tested on New German Credit, however, since its oracle produces a single, dataset-wide partition of concepts (Section~\ref{sec:oracle_tabular}) rather than a per-question grouping. Consequently, $\kappa(x, n_{\max})$ would be identical for every question in the dataset, leaving no variance in $\Delta\kappa(x)$ against which to correlate $\Delta CE(x)$.
\end{hypothesis_box}

\subsection{Evaluation of \ref{rq:convergence}: Baseline Equivalence}
The convergence of FELICS ($n_{\max}=1$) to the original \emph{Causal Concept Faithfulness} baseline is evaluated using an equivalence-testing rather than a difference-testing procedure, since the goal is to demonstrate the absence of a meaningful difference rather than the presence of one. Because the original metric does not support tabular data, this hypothesis is restricted to the BBQ and MedQA datasets.

\begin{hypothesis_box}
	For each dataset $X \in \{\text{BBQ}, \text{MedQA}\}$ and evaluated model $M$: restrict both the Matton et al.\ and the FELICS ($n_{\max}=1$) outputs to the set of questions present in both. For each shared question $x$, compute the mean CE across all its concepts---using $CE_{\mathrm{Matton}}(x,C_m)$ for the baseline and $CE_{\mathrm{adj}}(x,C_m,\{C_m\})$ for FELICS---as well as the question-level faithfulness score $F(x)$ obtained from each pipeline. Conduct a paired Two One-Sided Tests (TOST) procedure for both the CE sample and the faithfulness sample, using equivalence bounds of $\pm25\%$ and $\pm10\%$ of the standard deviation of the paired differences, respectively. The hypothesis is considered supported for a given dataset and model if both one-sided $p$-values of the TOST procedure fall below $0.20$, for both the CE comparison and the faithfulness comparison.
\end{hypothesis_box}

\subsection{Evaluation of \ref{rq:ranking}: Rank Preservation}
Rank preservation is evaluated on the BBQ, MedQA, and New German Credit datasets under both an atomic and a joint oracle configuration.

\begin{hypothesis_box}
	For each dataset $X \in \{\text{BBQ}, \text{MedQA}, \text{New German Credit}\}$, evaluated model $M$, and oracle configuration---$n_{\max}=1$, included as a sanity check since $CE_{\mathrm{raw}}$ and $CE_{\mathrm{adj}}$ coincide trivially for singleton groups, and the joint configuration ($n_{\max}=4$ for BBQ and MedQA, $n_{\max}=3$ for New German Credit)---collect, for every concept-level observation, the raw Shapley causal effect
	\begin{equation}
		CE_{\mathrm{raw}}(x,C_m,s) := \phi_m(v)
	\end{equation}
	obtained from Equation~\ref{eq:shapley} using the unadjusted value function $v$, alongside the adjusted Shapley causal effect $CE_{\mathrm{adj}}(x,C_m,s)$ obtained from Equation~\ref{eq:ce_adj}. Compute the Spearman rank correlation $\rho$ between both columns using a one-sided test (alternative: greater). The hypothesis is considered supported for a given dataset, model, and configuration if $\rho \ge 0.9$ and $p \le 0.05$.
\end{hypothesis_box}

\subsection{Evaluation of \ref{rq:faithfulness_increase}: Faithfulness Enhancement}
The effect of modeling interactions on faithfulness is evaluated on all datasets and models, comparing the atomic baseline against a joint configuration.

\begin{hypothesis_box}
	For each dataset $X \in \{\text{BBQ}, \text{MedQA}, \text{New German Credit}\}$ and evaluated model $M$: pair each question's atomic ($n_{\max}=1$) question-level faithfulness score with its joint-configuration counterpart ($n_{\max}=4$ for the natural-language datasets BBQ and MedQA, $n_{\max}=3$ for the tabular New German Credit dataset, consistent with the maximum group sizes established for the respective oracle implementations in Section~\ref{sec:joint_counterfactual}). Conduct a paired, one-sided Wilcoxon Signed-Rank Test with the alternative hypothesis that the joint-configuration faithfulness scores are stochastically greater than the atomic ones. The hypothesis is considered supported for a given dataset and model if the resulting $p$-value is $\le 0.05$.
\end{hypothesis_box}

\subsection{Evaluation of \ref{rq:motivating_example}: Motivating-Example Resolution}
The resolution of this specific case is verified through a targeted, qualitative analysis of the corresponding MedQA question, repeated for each evaluated model.

\begin{hypothesis_box}
	For each evaluated model $M$: run the identified MedQA question through both the original and the FELICS pipeline, and construct a table with columns \emph{Concept}, $CE_{\mathrm{original}}$, and $CE_{\mathrm{adj}}$, listing every concept present in the question. Isolate the two concepts of interest---\emph{the patient's eating disorder} and \emph{the patient's self-perception of weight}---and assert that $CE_{\mathrm{adj}} > CE_{\mathrm{original}}$ holds for both. Because this hypothesis concerns a single, specific question rather than a distribution of questions, no inferential statistical test is applicable; the analysis is instead reported as a qualitative case study, and the assertion is required to hold for every evaluated model.
\end{hypothesis_box}
